\documentclass[a4paper,11pt]{article}

\usepackage{fullpage}
\usepackage[utf8]{inputenc}
\usepackage[T1]{fontenc}
\usepackage{amsmath,amsthm}
\usepackage{amsfonts}
\usepackage{graphicx}
\usepackage{latexsym}
\usepackage{float}
\usepackage{comment}
\usepackage{array}
\usepackage{booktabs}

% --- Packages ajoutés pour l'enrichissement pédagogique ---
\usepackage{xcolor}
\usepackage{listings}
\usepackage[most]{tcolorbox}
\usepackage{hyperref}

%%%%%%%%%%%%%%% CONFIGURATION LISTINGS %%%%%%%%%%%%%%%%%%%

\definecolor{codegreen}{rgb}{0.25,0.49,0.25}
\definecolor{codegray}{rgb}{0.5,0.5,0.5}
\definecolor{codepurple}{rgb}{0.58,0,0.82}
\definecolor{codered}{rgb}{0.7,0.1,0.1}
\definecolor{codeblue}{rgb}{0.0,0.0,0.7}
\definecolor{backcolour}{rgb}{0.97,0.97,0.97}

\lstset{
    language=PHP,
    backgroundcolor=\color{backcolour},
    commentstyle=\color{codegreen}\itshape,
    keywordstyle=\color{codeblue}\bfseries,
    stringstyle=\color{codered},
    basicstyle=\ttfamily\small,
    breakatwhitespace=false,
    breaklines=true,
    captionpos=b,
    keepspaces=true,
    numbers=left,
    numbersep=8pt,
    numberstyle=\tiny\color{codegray},
    showspaces=false,
    showstringspaces=false,
    showtabs=false,
    tabsize=4,
    frame=single,
    rulecolor=\color{codegray},
    xleftmargin=1.5em,
    framexleftmargin=1.5em,
    literate={é}{{\'e}}1 {è}{{\`e}}1 {ê}{{\^e}}1 {à}{{\`a}}1 {ù}{{\`u}}1 {î}{{\^i}}1 {ô}{{\^o}}1
}

%%%%%%%%%%%%%%% ENVIRONNEMENTS TCOLORBOX %%%%%%%%%%%%%%%%%%

% Rappel de cours (fond bleu clair)
\newtcolorbox{rappel}[1][]{
    colback=blue!5!white,
    colframe=blue!50!black,
    fonttitle=\bfseries,
    title={\large Rappel de cours},
    #1
}

% Point clé (fond vert clair)
\newtcolorbox{pointcle}[1][]{
    colback=green!5!white,
    colframe=green!40!black,
    fonttitle=\bfseries,
    title={Point clé à retenir},
    #1
}

% Attention (fond orange clair)
\newtcolorbox{attention}[1][]{
    colback=orange!8!white,
    colframe=orange!60!black,
    fonttitle=\bfseries,
    title={Attention},
    #1
}

% Avertissement éthique (fond rouge)
\newtcolorbox{ethique}[1][]{
    colback=red!5!white,
    colframe=red!60!black,
    fonttitle=\bfseries\large,
    title={Avertissement éthique et légal},
    #1
}

%%%%%%%%%%%%%%% COMMANDES %%%%%%%%%%%%%%%%%%%%%%%%%%%

\newcommand{\itemb}{\item[$\bullet$]}

\newcommand{\Fd}{\mathbb{F}}
\newcommand{\Znat}{\mathbb{Z}}
\newcommand{\Nnat}{\mathbb{N}}

%%%%%%%%%%%%%%% ENVIRONEMENTS %%%%%%%%%%%%%%%%%%%%%%%

\theoremstyle{plain}
\newtheorem{lemme}{Lemme}
\newtheorem{algorithme}{Algorithme}
\newtheorem{corolaire}{Corollaire}
\newtheorem{propriete}{Propri\'et\'e}
\newtheorem{proposition}{Proposition}
\newtheorem{theoreme}{Th\'eor\`eme}
\newtheorem{definition}{D\'efinition}

\theoremstyle{definition}
\newtheorem{exercice}{Exercice}
\newtheorem{remarque}{Remarque}
\newtheorem{exemple}{Exemple}

\hypersetup{
    colorlinks=true,
    linkcolor=blue!60!black,
    urlcolor=blue!60!black
}

%%%%%%%%%%%%%% DOCUMENTS %%%%%%%%%%%%%%%%%%%%%%%%%%%%


\begin{document}

\hbox{\raisebox{0.4em}{\vrule depth 0pt height 0.5pt width \textwidth}}
\noindent \textbf{Université de Perpignan} \hfill  L3 Informatique \\
 \hfill C. Legrand \hfill  Année \the\year \\
\smallskip
\begin{center}
{
  {\scshape \large TD7 Sécurité}  \\
 Injection de code
}
\end{center}
\hbox{\raisebox{0.4em}{\vrule depth 0pt height 0.9pt width \textwidth}}

\bigskip

% --- Avertissement éthique ---
\begin{ethique}
Les techniques présentées dans ce TD sont étudiées dans un \textbf{cadre strictement académique} afin de comprendre les mécanismes utilisés par les attaquants et ainsi pouvoir mieux s'en protéger.

\textbf{Toute utilisation de ces techniques en dehors du cadre pédagogique est illégale} et passible de sanctions pénales (chapitre III du Code pénal, articles 323-1 à 323-8 -- \textit{Des atteintes aux systèmes de traitement automatisé de données}) :
\begin{itemize}
    \item Accès ou maintien frauduleux dans un STAD (art.~323-1) : \textbf{3 ans} d'emprisonnement et \textbf{100\,000\,\texteuro{}} d'amende
    \item Entrave au fonctionnement d'un STAD (art.~323-2) : \textbf{5 ans} d'emprisonnement et \textbf{150\,000\,\texteuro{}} d'amende
    \item Introduction frauduleuse de données dans un STAD (art.~323-3) : \textbf{5 ans} d'emprisonnement et \textbf{150\,000\,\texteuro{}} d'amende
\end{itemize}
Peines portées à \textbf{7 ans} et \textbf{300\,000\,\texteuro{}} lorsque le système visé traite des données à caractère personnel mis en \oe uvre par l'État, et jusqu'à \textbf{10 ans} et \textbf{300\,000\,\texteuro{}} en bande organisée (art.~323-4-1). \textit{Peines mises à jour par la loi LOPMI n\textsuperscript{o}2023-22 du 24 janvier 2023.}
\end{ethique}

\bigskip

\begin{tcolorbox}[colback=gray!5!white, colframe=gray!60!black, title={Mise en place avec Docker}]
Pour ce TD, l'environnement (serveur Apache + PHP + MySQL) est fourni sous forme de conteneurs Docker. Aucune machine virtuelle n'est nécessaire.

\medskip
\textbf{Prérequis : installer Docker}

\medskip
\textit{Sous Ubuntu/Debian :}
\begin{lstlisting}[language=bash,numbers=none]
# Installation de Docker
sudo apt update
sudo apt install -y docker.io docker-compose-v2
# Ajouter votre utilisateur au groupe docker
sudo usermod -aG docker $USER
# Deconnectez-vous puis reconnectez-vous pour
# que le changement de groupe prenne effet
\end{lstlisting}

\textit{Sous macOS/Windows :} installer \href{https://www.docker.com/products/docker-desktop/}{Docker Desktop}.

\medskip
\textbf{Lancement de l'environnement}

\begin{lstlisting}[language=bash,numbers=none]
# Se placer dans le repertoire du TD
cd travaux-diriges/04-injection-code/td07/

# Lancer les conteneurs (premiere fois : ~2 min)
docker compose up -d

# Verifier que tout fonctionne
docker compose ps
\end{lstlisting}

Au premier lancement, Docker :
\begin{itemize}
    \item construit un conteneur Apache+PHP avec l'extension \texttt{mysqli},
    \item lance un conteneur MySQL 8.0 avec la base \texttt{l2info},
    \item charge automatiquement les tables et données depuis \texttt{bd.sql}.
\end{itemize}

\medskip
\textbf{Accès aux pages web}

Les fichiers du TD sont accessibles dans le navigateur à l'adresse :

\begin{center}
\texttt{http://localhost:8080/td7/}
\end{center}

Par exemple : \texttt{http://localhost:8080/td7/page-5-commentaires.php}

\medskip
\textbf{Arrêt et nettoyage}

\begin{lstlisting}[language=bash,numbers=none]
# Arreter les conteneurs (les donnees sont conservees)
docker compose stop

# Relancer
docker compose start

# Arreter et supprimer tout (conteneurs + donnees)
docker compose down -v
\end{lstlisting}

\medskip
\begin{attention}
Dans les fichiers PHP, le serveur MySQL est accessible via le nom d'hôte \texttt{"db"} (le nom du service Docker), et non \texttt{"localhost"}. C'est le fonctionnement normal de Docker : chaque service est un conteneur séparé, identifié par son nom.
\end{attention}

\end{tcolorbox}

\bigskip


%%%%%%%%%%%%%%%%%%%%%%%%%%%%%%%%%%%%%%%%%%%%%%%%%%%%%%%%%%%%%%%%%%
%% EXERCICE 1 : Manipulation de formulaire HTML
%%%%%%%%%%%%%%%%%%%%%%%%%%%%%%%%%%%%%%%%%%%%%%%%%%%%%%%%%%%%%%%%%%

\begin{exercice}[Manipulation de formulaire HTML]

On considère une application web d'achat en ligne dont le formulaire HTML est le suivant (fichier \texttt{achat-iphone5.html}) :

\begin{lstlisting}[language=HTML]
<html>
<body>

<form method="post" action="enregistre-achat.php">
Produit: iphone5 <br/>
Prix : 500 euros <br/>
Quantit&eacute <input type="text" name ="quantite">
  (Quantite max est 90) <br/>
<input type="hidden" name="prix" value="500" >
<input type="submit" value="buy">
</form>

</body>
</html>
\end{lstlisting}

Le script PHP \texttt{enregistre-achat.php} enregistre l'achat en base de données en utilisant les valeurs envoyées par le formulaire.

\begin{enumerate}
    \item Analysez le formulaire HTML. Identifiez la vulnérabilité liée à la manière dont le prix est transmis au serveur.
    \item Proposez une méthode pour acheter un iPhone au prix de 5\,\texteuro{} au lieu de 500\,\texteuro{}.
    \item Proposez une correction côté serveur pour empêcher cette attaque.
\end{enumerate}

\begin{pointcle}
Les champs \texttt{hidden} et toute donnée côté client peuvent être modifiés par l'utilisateur. La validation doit toujours se faire côté serveur.
\end{pointcle}

\end{exercice}


%%%%%%%%%%%%%%%%%%%%%%%%%%%%%%%%%%%%%%%%%%%%%%%%%%%%%%%%%%%%%%%%%%
%% EXERCICE 2 : Cross-Site Scripting (XSS)
%%%%%%%%%%%%%%%%%%%%%%%%%%%%%%%%%%%%%%%%%%%%%%%%%%%%%%%%%%%%%%%%%%

\begin{exercice}[Cross-Site Scripting (XSS)]

\begin{rappel}
Il existe trois types de XSS :
\begin{itemize}
    \item \textbf{XSS réfléchi (reflected)} : le payload est inclus dans la requête (URL, paramètre GET/POST) et renvoyé immédiatement dans la réponse sans être stocké.
    \item \textbf{XSS stocké (stored)} : le payload est enregistré sur le serveur (base de données, fichier) et affiché à tous les visiteurs ultérieurs de la page.
    \item \textbf{XSS basé sur le DOM (DOM-based)} : le payload est interprété directement par le JavaScript côté client, sans transiter par le serveur.
\end{itemize}
\end{rappel}

\medskip

\textbf{a) XSS stocké -- alerte JavaScript}

\medskip

On considère la page \texttt{page-5-commentaires.php} dont le code est le suivant :

\begin{lstlisting}[language=PHP]
<html>
<H1> Computer Security  </H1>
<IMG src="img_security.jpeg"> Computer security...
<BR/><hr/>
<h3> Visitors comments :</h3>
<?php
$conn = new mysqli("db", "l2info", "l2info", "l2info");
if ($conn->connect_error) {
  die("Connection failed: " . $conn->connect_error);
}

if( isset($_POST['login']) && isset($_POST['commentaire']) )
{
    $sql = 'INSERT INTO Commentaires(login,commentaire)
      VALUES("' . $_POST['login']. '","'
      . $_POST['commentaire'].'");';
    $result0 = $conn->query($sql);
}

$sql = 'SELECT id,login, commentaire FROM Commentaires
  ORDER BY id DESC LIMIT 4 ;';
$result = $conn->query($sql);

if ($result->num_rows > 0) {
    echo "<ul>";
    while($row = $result->fetch_assoc()) {
        echo "<li>";
        echo "[" . $row["login"]. "]: "
          . $row["commentaire"]."<br>";
        echo "</li>";
    }
    echo "</ul>";
} else {
    echo "Aucun commentaire.";
}
$conn->close();
?>
<BR/><hr/><BR/>
<form name="postcomment" method="POST"
  action="page-5-commentaires.php">
Pseudo : <input name="login"> <BR/>
Commentaire : <input type="text" name="commentaire"><BR/>
<input type="submit">
</form>
</html>
\end{lstlisting}

\begin{enumerate}
    \item Expliquez ce que fait cette page.
    \item Proposez une injection dans le champ commentaire pour qu'un message d'alerte JavaScript s'affiche au chargement de la page.
\end{enumerate}

\medskip

\textbf{b) XSS stocké -- défacement}

\medskip

En utilisant la même page, injectez du HTML frauduleux pour faire apparaître une fausse section \og Références \fg{} avec un lien vers un site fictif.

\medskip

\textbf{c) XSS réfléchi via GET}

\medskip

On considère le code PHP suivant :

\begin{lstlisting}[language=PHP]
<html>
<body>
<?php
if( isset($_GET['UserInput'])){
  $out ='vous avez saisi: "'
    . $_GET['UserInput'] . '".';
}
else {
  $out = '<form method="GET">
    saisissez quelque chose ici :';
  $out .= '<input name="UserInput" size="50">';
  $out .= '<input type="submit">';
  $out .= '</form>';
}
print $out;
?>
</body>
</html>
\end{lstlisting}

Soumettez une entrée qui fait apparaître une alerte JavaScript dans le résultat.

\medskip

\textbf{d) XSS réfléchi via POST}

\medskip

On considère maintenant le code suivant :

\begin{lstlisting}[language=PHP]
<html>
<body>
<?php
if( isset($_POST['UserInput'])){
  $out ='vous avez saisi: "'
    . $_POST['UserInput'] . '".';
}
else {
  $out='Aucune donnee transmise au script';
}
print $out;
?>
</body>
</html>
\end{lstlisting}

Construisez une page HTML externe contenant un formulaire POST caché qui soumet automatiquement un payload XSS à cette page.

\medskip

\textbf{e) XSS basé sur le DOM}

\medskip

On considère la page HTML suivante :

\begin{lstlisting}[language=HTML]
<html>
<body>
<div id="welcome"></div>
<script>
  var name = document.location.hash.substring(1);
  document.getElementById("welcome").innerHTML
    = "Bienvenue " + name;
</script>
</body>
</html>
\end{lstlisting}

\begin{enumerate}
    \item Construisez une URL qui exploite cette vulnérabilité pour exécuter du JavaScript.
    \item Expliquez pourquoi ce type de XSS est plus difficile à détecter côté serveur.
\end{enumerate}

\begin{attention}
Le XSS stocké est le plus dangereux car il affecte tous les visiteurs de la page, sans nécessiter de lien piégé.
\end{attention}

\end{exercice}


%%%%%%%%%%%%%%%%%%%%%%%%%%%%%%%%%%%%%%%%%%%%%%%%%%%%%%%%%%%%%%%%%%
%% EXERCICE 3 : Injection SQL
%%%%%%%%%%%%%%%%%%%%%%%%%%%%%%%%%%%%%%%%%%%%%%%%%%%%%%%%%%%%%%%%%%

\begin{exercice}[Injection SQL]

\begin{rappel}
\textbf{Différence entre XSS et injection SQL :}
\begin{itemize}
    \item Le \textbf{XSS} injecte du code malveillant dans le \textit{navigateur de la victime} (côté client). L'attaque vise les utilisateurs du site.
    \item L'\textbf{injection SQL} injecte du code malveillant dans la \textit{base de données du serveur} (côté serveur). L'attaque vise directement les données du site.
\end{itemize}
\end{rappel}

\medskip

\textbf{a) Analyse du code vulnérable}

\medskip

On considère la page \texttt{page-transaction-sql.php} :

\begin{lstlisting}[language=PHP]
<html>
<?php
$servername = "db";
$username = "l2info";
$password = "l2info";
$dbname = "l2info";

$conn = new mysqli($servername, $username,
  $password, $dbname);
if ($conn->connect_error) {
  die("Connection failed: "
    . $conn->connect_error);
}
echo 'Nom : "' . $_POST['UserInput'].'"<BR/>';
$sql = 'SELECT montant FROM Transaction WHERE nom="'
  . $_POST['UserInput']  . '";';
echo "Requete SQL=" . $sql . "<BR/>";
$result = $conn->query($sql);
echo "Resultat:<BR/>";
if ($result->num_rows > 0) {
  while($row = $result->fetch_assoc()) {
    echo "Nom: ". $_POST['UserInput']
      ." Montant :" . $row["montant"]."<br/>";
  }
} else {
  echo "0 results";
}
$conn->close();
?>
</html>
\end{lstlisting}

Et le formulaire \texttt{injection-sql.html} :

\begin{lstlisting}[language=HTML]
<html>
<body>
<form name="exploit-sql" method="POST"
  action="http://localhost:8080/td7/page-transaction-sql.php">
<input name="UserInput">
<input type="submit">
</form>
</body>
</html>
\end{lstlisting}

\begin{enumerate}
    \item Expliquez comment la requête SQL est construite. Identifiez la vulnérabilité.
\end{enumerate}

\medskip

\textbf{b) Afficher toute la table}

\medskip

Proposez une injection pour afficher l'intégralité de la table \texttt{Transaction}. Expliquez pas-à-pas la requête résultante.

\medskip

\textbf{c) Tentative de DROP TABLE}

\medskip

Tentez de supprimer la table \texttt{Test} par injection SQL. Expliquez pourquoi cela peut échouer avec MySQL.

\medskip

\textbf{d) UNION-based}

\medskip

Extraire les données de la table \texttt{Test} (colonne \texttt{val}) via une injection \texttt{UNION SELECT}.

\textit{Indice :} commencez par déterminer le nombre de colonnes du \texttt{SELECT} original en utilisant \texttt{ORDER BY 1}, \texttt{ORDER BY 2}, etc. jusqu'à obtenir une erreur. Puis construisez le \texttt{UNION SELECT} avec le bon nombre de colonnes.

\medskip

\textbf{e) Réflexion sur le blind SQL}

\medskip

On suppose maintenant que la page n'affiche plus les résultats de la requête mais seulement \og Transaction trouvée \fg{} ou \og Aucune transaction \fg{}.

\begin{enumerate}
    \item Comment pourrait-on déterminer si la version de MySQL est 5.x ? (concept \textit{boolean-based blind SQL injection})
    \item Quel serait le principe pour extraire des données caractère par caractère ?
\end{enumerate}

\begin{pointcle}
L'injection SQL fonctionne car la requête mélange code SQL et données utilisateur dans une même chaîne. La solution est de les séparer avec des \textbf{requêtes préparées} (\textit{prepared statements}).
\end{pointcle}

\end{exercice}


%%%%%%%%%%%%%%%%%%%%%%%%%%%%%%%%%%%%%%%%%%%%%%%%%%%%%%%%%%%%%%%%%%
%% EXERCICE 4 : Contre-mesures
%%%%%%%%%%%%%%%%%%%%%%%%%%%%%%%%%%%%%%%%%%%%%%%%%%%%%%%%%%%%%%%%%%

\begin{exercice}[Contre-mesures]

\textbf{a) Corriger le XSS}

\medskip

Modifiez le code de \texttt{page-5-commentaires.php} pour empêcher l'injection XSS. Quelle fonction PHP permet d'échapper les caractères spéciaux HTML ? Donnez le code corrigé de la ligne d'affichage des commentaires.

\medskip

\textbf{b) Corriger l'injection SQL}

\medskip

Réécrivez \texttt{page-transaction-sql.php} en utilisant des requêtes préparées PDO au lieu de la concaténation avec \texttt{mysqli}. Donnez le code complet corrigé.

\medskip

\textbf{c) Défenses complémentaires}

\begin{enumerate}
    \item Qu'est-ce que le header HTTP \texttt{Content-Security-Policy} et comment protège-t-il contre le XSS ?
    \item Qu'est-ce qu'un cookie \texttt{HttpOnly} et pourquoi est-ce utile ?
    \item Pourquoi le principe du moindre privilège est-il important pour la base de données ?
\end{enumerate}

\begin{pointcle}
La défense en profondeur combine plusieurs mécanismes : échappement en sortie, requêtes préparées, CSP, cookies \texttt{HttpOnly}, validation des entrées.
\end{pointcle}

\end{exercice}


\end{document}
